% ===================================================================
%  BAGAN No. 1 — Sekolah, sekolah menengah, ruang kelas.
%
%  Ceci n'est PAS une transcription : c'est une traduction, faite en
%  2026, en indonesien d'aujourd'hui. D'ou trois differences avec
%  texto/io et texto/fr, et trois seulement :
%
%   * pas de \nl ni de \cc — la traduction ne suit aucune ligne
%     imprimee, et les lignes de ce fichier ne sont que du confort ;
%   * pas de \begin{VUpage} — elle n'a ni feuillet ni folio ;
%   * le gras se pose sur le TERME seul, ponctuation dehors.
%
%  Les cles « %%K » sont celles de texto/io, macro pour macro pour
%  l'apparat, et les renvois \textsuperscript{(n)} visent les MEMES
%  objets de la planche, dans le MEME ordre.
%
%  CETTE COLONNE N'A AUCUNE VOISINE DANS LE RELEVE, ET C'EST
%  JUSTEMENT SA DIFFICULTE. Les colonnes defendues jusqu'ici l'etaient
%  contre une langue presente ailleurs dans texto/ — le cantonais
%  contre le mandarin, le marathe contre le hindi, l'ourdou contre le
%  chahmoukhi et le hindi a la fois. L'indonesien a bien une voisine,
%  et de tres pres : le malais de Malaisie est la MEME langue a un
%  standard pres. Mais il n'est dans aucune colonne, rien ne le
%  signalerait a l'oeil, et il n'y a donc rien a comparer. Le
%  controle doit se faire seul.
%
%  kolonoj.py LE TIENT PAR LE LEXIQUE, ET SUR DEUX FRONTS.
%
%   * LE PREMIER EST GEOGRAPHIQUE. Les deux standards se separent sur
%     une liste courte et tres connue de mots quotidiens : cikgu
%     contre guru, tandas contre toilet, stesen contre stasiun, ubat
%     contre obat, wang contre uang, beg contre tas, kerusi contre
%     kursi, almari contre lemari, tingkap contre jendela, bomba
%     contre pemadam kebakaran, tarikh contre tanggal, tayar contre
%     ban, basikal contre sepeda, hospital contre rumah sakit. Ce
%     sont ceux-la qu'on releve.
%
%   * LE SECOND EST UNE DATE, ET IL EST PROPRE A CE LIVRET.
%     L'orthographe d'avant la reforme de 1972 — « boekoe » pour
%     buku, « djalan » pour jalan, « tjelana » pour celana, « jang »
%     pour yang — est exactement celle qu'un imprimeur aurait employee
%     en 1926, l'annee du fac-simile. C'est le seul piege du releve
%     ou la faute serait CONTEMPORAINE de la source : une colonne
%     ecrite « comme a l'epoque » serait plus fidele a la date et
%     moins fidele a la consigne, qui est d'ecrire l'indonesien de
%     2026.
%
%  ET QUATRE MOTS N'ONT PAS ETE RELEVES, EXPRES, PARCE QU'ILS SONT DES
%  FAUX AMIS. « pejabat » est le bureau en Malaisie et le
%  fonctionnaire en Indonesie ; « budak » est l'enfant la-bas et
%  l'esclave ici ; « polis » est la police la-bas et la police
%  d'assurance ici — et le tableau 7 met justement en scene un agent
%  d'assurances ; « kedai » et « lori » existent des deux cotes avec
%  des emplois differents. Un controle qui crie sur un mot juste finit
%  desarme : c'est la lecon de « دی » et « دے » a la colonne ourdoue.
%
%  LE VASISTAS (78) EST NOMME PAR UN MOT NEERLANDAIS, ET C'EST LA
%  MEILLEURE RENCONTRE DU RELEVE. Le francais nomme cette fenetre par
%  une question allemande — « was ist das » —, l'indonesien la nomme
%  par un mot neerlandais, « boven », de bovenlicht : le jour d'en
%  haut. Deux langues, deux voisins du nord, et la meme facon
%  d'emprunter au voisin le nom d'une chose qu'on a chez soi. La serie
%  du vasistas continue donc — un seul mot, quarante fois sur
%  quarante et une.
%
%  L'INDONESIEN EST A TETE INITIALE : le determinant suit le
%  determine, « ruang kelas » et non l'inverse. parigi.py ne compte
%  donc que huit paires pour ce tableau, contre trente-cinq au tableau
%  9 ourdou : la regle modificateur-tete, qui a servi cinq familles a
%  tete finale, ne mord presque pas ici. C'est la premiere colonne du
%  releve depuis longtemps ou l'ordre des renvois se suit sans
%  parade.
%
%  LE VOCABULAIRE DE LA CLASSE EST PARTAGE ENTRE TROIS COUCHES, ET LA
%  COLONNE LES GARDE TOUTES LES TROIS :
%
%    papan tulis (1)  le tableau noir   sabak (42)    l'ardoise
%    penghapus (2)    le torchon        kapur (3)     la craie
%    lemari (15)      l'armoire         kamus (30)    le dictionnaire
%    buku tulis (28)  le cahier         jangka (26)   le compas
%    penggaris (50)   la regle          pisau raut (56) le canif
%    tas sekolah (57) le cartable       tungku (46)   le poele
%    halaman (75)     la cour           boven (78)    le vasistas
%
%  Malais de fond pour les objets qu'on manie, neerlandais pour ce que
%  l'ecole coloniale a apporte (grip, spons, boven), arabe et sanskrit
%  pour l'abstrait (guru, ilmu, sejarah, kertas). Aucune ne se sent
%  etrangere en 2026, et c'est ce qui distingue cette colonne de
%  toutes les precedentes : ici l'emprunt n'est pas une reponse a un
%  trou, c'est le tissu meme de la langue.
%
%  CORRECTION, ECRITE APRES LE TABLEAU 5 : IL Y A QUATRE COUCHES, ET
%  NON TROIS. Il manquait le HOKKIEN, apporte par les marchands et les
%  artisans chinois — « cat » la peinture, de 漆 ; « loteng » le
%  grenier, de 樓頂 ; et plus loin dans le volume « teko », « bakmi »,
%  « tauge ». Le tableau 5, qui est celui des metiers du batiment, l'a
%  fait paraitre d'un coup. La faute n'etait pas dans le compte mais
%  dans la methode : on avait annonce le partage sur un seul tableau
%  au lieu d'attendre que le volume le montre. Une couche entiere peut
%  se cacher derriere quatre tableaux.
% ===================================================================

%%K t01-apar-0 apar
\VUsaut{2.0mm}
\VUcentre{12.6pt}{120}{BUKU PENJELASAN}
\VUsaut{3.2mm}
\VUcentre{8.4pt}{160}{UNTUK}
\VUsaut{2.6mm}
\VUtitre{15.6pt}{0}{Bagan Bantu Delmas}
\VUsaut{3.4mm}
\VUornamento{9pt}
\VUsaut{5.0mm}
\VUcentre{13.2pt}{90}{SERI PERTAMA}
\VUsaut{3.0mm}
\VUfilet{16mm}
\VUsaut{5.6mm}

%%K t01-apar-1 apar
\VUtitre{15.0pt}{60}{BAGAN No. 1}
\VUsaut{1.4mm}
\VUtitre{11.4pt}{0}{Sekolah, sekolah menengah, ruang kelas.}
\VUsaut{2.2mm}
\VUfilet{20mm}
\VUsaut{6.4mm}

%%K t01-tit-1 sub
\VUpk{10.2pt}{40}{Gambaran umum.}
\VUsaut{3.0mm}

%%K t01-01-1 p
1. --- Bagan ini menggambarkan sebuah \VUgras{ruang kelas} di sebuah
\VUgras{sekolah menengah}. Di ruang kelas ini ada delapan
\VUgras{meja} \textsuperscript{(18)} dan delapan \VUgras{bangku}
\textsuperscript{(32)}. Di setiap bangku duduk tiga murid.
\VUgras{Guru} \textsuperscript{(7)} duduk di \VUgras{kursi}
\textsuperscript{(8)} di \VUgras{meja gurunya} \textsuperscript{(6)}.

%%K t01-02-1 p
2. --- Di sebelah kiri meja guru berdiri \VUgras{lemari}
\textsuperscript{(15)} \VUgras{berkaca}. Di sebelah kanan ada
\VUgras{papan tulis} \textsuperscript{(1)} dengan \VUgras{penghapusnya}
\textsuperscript{(2)} dan \VUgras{kapurnya} \textsuperscript{(3)}.

%%K t01-02-2 p
Di atas meja guru tergantung \VUgras{papan peraga} \textsuperscript{(9,
11, 12)} dan sebuah \VUgras{peta} \textsuperscript{(10)}.

%%K t01-03-1 p
3. --- Tidak semua murid duduk di \VUgras{tempatnya}
\textsuperscript{(17)}. Ioannes \textsuperscript{(4)} berdiri di depan
papan tulis ; ia memegang \VUgras{penghapus} itu dan menghapus
\VUgras{gambar-gambar geometri.}

%%K t01-03-2 p
Petrus berdiri dekat meja guru ; ia mengumpulkan \VUgras{tugas
menulis} \textsuperscript{(5)} dan membawanya kepada guru.

%%K t01-04-1 p
4. --- Guru \VUgras{ini} adalah guru \VUgras{bahasa-bahasa hidup.} Ia
mengajar murid-murid berbicara, membaca dan menulis bahasa asing
\VUgras{(Jerman, Inggris, Spanyol, Italia, Rusia atau Prancis)}.

%%K t01-05-1 p
5. --- Ruang kelas itu sangat luas dan sangat terang. Di ujung ruang
ada \VUgras{jendela} lebar dengan dua \VUgras{boven}
\textsuperscript{(78)} dan sebuah \VUgras{pintu berkaca} untuk
mengudarakan dan menerangi ruang itu.

%%K t01-05-2 p
Ruang kelas ini tidak dingin. Di tengahnya, untuk menghangatkannya,
ada \VUgras{tungku} \textsuperscript{(46)} yang sangat baik dengan
\VUgras{cerobongnya} \textsuperscript{(45)}.

%%K t01-05-3 p
Pada dinding terpasang \VUgras{gantungan baju} \textsuperscript{(58)} ;
di situ tergantung topi dan mantel murid-murid.

%%K t01-tit-2 sub
\VUsaut{5.2mm}
\VUpk{10.2pt}{40}{Halaman bermain}
\VUsaut{3.6mm}

%%K t01-06-1 p
6. --- \VUgras{Jam dinding} \textsuperscript{(63)} menunjukkan pukul
sembilan lewat lima menit.

%%K t01-06-2 p
\VUgras{Waktu istirahat} \textsuperscript{(76)} baru saja berakhir dan
pelajaran dimulai lagi. Tempat bermain terletak di \VUgras{halaman}
\textsuperscript{(75)}. Kita melihat beberapa murid yang sedang
bermain. Seorang \VUgras{pengawas} \textsuperscript{(74)} menjaga
mereka.

%%K t01-07-1 p
7. --- Di halaman kita juga melihat beberapa orang lain, yaitu
\VUgras{penilik sekolah} \textsuperscript{(73)}, \VUgras{kepala
sekolah} \textsuperscript{(71)} dan \VUgras{wakil kepala sekolah}
\textsuperscript{(72)}.

%%K t01-07-2 p
Penilik memeriksa sekolah-sekolah, kepala sekolah memimpin sekolah
menengah itu, dan wakilnya mengawasi pelajaran.

%%K t01-07-3 p
Orang keempat adalah \VUgras{penjaga pintu} \textsuperscript{(77)}. Ia
membawa buku daftar di bawah lengannya untuk mencatat yang tidak
hadir.

%%K t01-08-1 p
8. --- Di ujung halaman ada \VUgras{alat-alat senam}
\textsuperscript{(70)} dan bengkel untuk \VUgras{kerajinan tangan}
\textsuperscript{(69)}.

%%K t01-08-2 p
Di sebelah kanan bagan ini kita mulai melihat \VUgras{ruang-ruang
kelas} untuk \VUgras{musik} dan \VUgras{menggambar}.

%%K t01-noto-1 noto
\VUnotes{143.2mm}{%
(*) Untuk \textit{nama baptis} pun disarankan menyalinnya menurut
bentuk Latin. Itulah satu-satunya jalan untuk menghindari perbedaan
yang tak terhitung. Lagi pula, bagaimana memilih di antara
\textit{Giovanni, Jean, Johann, Jan, Hans, John, Ivan,} dan
sebagainya? Apakah kita akan membuat kamus khusus untuk nama baptis?
Marilah kita ambil bentuk nominatifnya dari bahasa Latin dan
mengatakan : \VUgras{Ioannes, Ioanna; Iulius, Iulia; Iakobus;
Andreas; Lukas.} (Gram. Kompleta).%
}

%%K t01-tit-3 sub
\VUpk{10.2pt}{40}{Gambaran terperinci.}
\VUsaut{2.4mm}

%%K t01-tit-4 sub
\VUpk{10.2pt}{40}{Murid-murid yang baik.}
\VUsaut{3.0mm}

%%K t01-09-1 p
9. --- Murid-murid di bangku pertama di sebelah kanan meja guru
adalah \VUgras{Henrikus} (*), \VUgras{Stefanus} dan \VUgras{Karolus.}

%%K t01-09-2 p
Henrikus sangat penuh perhatian dan rajin ; ia mendengarkan guru.

Di atas mejanya ada \VUgras{buku tulis} \textsuperscript{(28)},
\VUgras{buku} \textsuperscript{(29)}, \VUgras{kamus}
\textsuperscript{(30)} dan \VUgras{kotak pensil} \textsuperscript{(31)}.
Bukunya mempunyai banyak \VUgras{halaman} ; setiap bab berisi beberapa
\VUgras{paragraf} dan setiap \VUgras{baris} banyak \VUgras{kata.}

%%K t01-09-4 p
Tetangganya, Stefanus \textsuperscript{(24)}, sangat bersemangat. Ia
mencatat di buku tulisnya pelajaran untuk pertemuan berikutnya.
\VUgras{Tulisannya} bagus. Bukunya tertutup ; kita hanya melihat
\VUgras{sampulnya} \textsuperscript{(27)}. Di dekatnya ada
\VUgras{jangkanya} \textsuperscript{(26)}.

%%K t01-09-5 p
Karolus \textsuperscript{(23)} memegang bukunya dengan tangan ; ia
membukanya untuk membaca ulang pelajarannya. Seperti setiap murid, ia
mempunyai \VUgras{tempat tinta} \textsuperscript{(25)} di depannya. Ia
penurut.

%%K t01-10-1 p
10. --- Di meja sebelahnya ada \VUgras{Ludovikus,} \VUgras{Antonius}
dan \VUgras{Paulus.} Ludovikus menghafalkan \VUgras{pelajarannya}
\textsuperscript{(22)} dan menjawab \VUgras{pertanyaan-pertanyaan}
guru. Antonius \textsuperscript{(21)} sangat tidak memperhatikan ;
kadang-kadang guru bahkan menghukumnya. Paulus
\textsuperscript{(20)} kawan yang baik, ramah dan suka menolong. Ia
sangat penuh perhatian.

%%K t01-tit-5 sub
\VUsaut{4.2mm}
\VUpk{10.2pt}{40}{Murid-murid yang buruk.}
\VUsaut{3.2mm}

%%K t01-11-1 p
11. --- Di belakang Henrikus duduk \VUgras{Georgius}
\textsuperscript{(38)}. Ia bukan anak yang jahat, tetapi ia
\VUgras{suka mengobrol,} tidak memperhatikan dan nakal.
\VUgras{Keisengannya} dan \VUgras{kebandelannya} menimbulkan
\VUgras{keributan} selama pelajaran, dan sering ia pantas mendapat
\VUgras{teguran} keras. \VUgras{Buku buramnya} \textsuperscript{(35)}
kotor ; ia \VUgras{menodainya dengan tinta} dari \VUgras{penanya}
\textsuperscript{(34)}, dan karena itu ia selalu memakai \VUgras{karet
penghapusnya} \textsuperscript{(36)} ; \VUgras{map bukunya}
\textsuperscript{(33)} sobek, sebab ia sangat ceroboh. Mengapa pula ia
membawa \VUgras{tempat pensilnya} \textsuperscript{(37)} ke ruang kelas
\VUgras{bahasa-bahasa hidup?}

%%K t01-11-2 p
Tetangganya, Edmundus \textsuperscript{(39)}, tidak menyukainya.
Memang ia agak lamban, tetapi ia pendiam dan tenang. Ia tidak
mengobrol ; hanya, alih-alih mendengarkan, ia lebih suka membaca
\VUgras{cerita-cerita} dalam \VUgras{buku bacaannya.}

%%K t01-11-3 p
Murid ketiga di bangku ini adalah \VUgras{Ludovikus}
\textsuperscript{(40)}. Ia tekun. Ia menulis di atas \VUgras{lembar
kertas} putih. Di depannya ia meletakkan \VUgras{kertas isapnya}
\textsuperscript{(41)}.

%%K t01-12-1 p
12. --- Murid-murid di bangku kedua, di sebelah kiri guru, masih
sangat muda. Yang pertama, si kecil \VUgras{Emilius,} belum pandai
benar menulis ; ia membawa \VUgras{sabaknya} \textsuperscript{(42)},
\VUgras{pensil sabaknya} dan \VUgras{sponsnya} \textsuperscript{(43)}.

%%K t01-12-2 p
Di sebelahnya ada \VUgras{Augustus} dan \VUgras{Bernardus}
\textsuperscript{(44)}. Yang terakhir ini bekerja dengan baik, tetapi
\VUgras{pakaiannya} sangat tidak terurus.

%%K t01-13-1 p
13. --- Di belakang mereka duduk tiga anak \VUgras{pemalas.
Albertus} tidak menghafal pelajarannya. \VUgras{Iakobus}
\textsuperscript{(47)} tidur alih-alih mendengarkan.
\VUgras{Kemalasannya} mendatangkan \VUgras{teguran} dan bahkan
\VUgras{hukuman} baginya. Akhirnya \VUgras{Kristianus}
\textsuperscript{(48)} terlalu tidak bersungguh-sungguh. Ia
\VUgras{berkelakuan} buruk dan sama sekali tidak berusaha
memperbaiki diri.

%%K t01-14-1 p
14. --- Sebaliknya, tetangga-tetangganya berkelakuan baik.
\VUgras{Ernestus} \textsuperscript{(51)} menyukai ketertiban dan
keteraturan ; ia selalu memakai \VUgras{penggarisnya}
\textsuperscript{(50)} dan \VUgras{pensilnya} \textsuperscript{(49)}.
Ia \VUgras{murid luar.}

%%K t01-14-2 p
\VUgras{Franciskus} \textsuperscript{(52)} \VUgras{murid asrama} ; ia
berpakaian seragam, makan dan tidur di sekolah menengah itu. Ia
\VUgras{kawan} yang baik.

%%K t01-14-3 p
Mauricius meraut \VUgras{pensilnya} dengan \VUgras{pisau rautnya}
\textsuperscript{(56)} ; ia sangat cermat.

%%K t01-14-4 p
Ia membawa semua bukunya, \VUgras{tata bahasanya}
\textsuperscript{(55)}, \VUgras{buku catatannya} \textsuperscript{(54)},
dan bahkan \VUgras{alas tulisnya} \textsuperscript{(53)}. \VUgras{Tas
sekolahnya} \textsuperscript{(57)} ada di lantai di belakangnya.

%%K t01-14-5 p
Kedua meja di ujung ruang kelas ditempati oleh \VUgras{murid-murid
yang baik} \textsuperscript{(19)}.

%%K t01-tit-6 sub
\VUsaut{4.6mm}
\VUpk{10.2pt}{40}{Menggambar dan musik.}
\VUsaut{3.4mm}

%%K t01-15-1 p
15. --- Di \VUgras{ruang menggambar} ada \VUgras{model-model} yang
indah tergantung di dinding dan sebuah \VUgras{lampu}
\textsuperscript{(62)} dengan \VUgras{kap lampunya,} tergantung di
langit-langit. \VUgras{Gurunya} \textsuperscript{(60)} sangat sabar ;
ia membetulkan \VUgras{gambar-gambar} \textsuperscript{(61)} murid dan
mengajar mereka menggambar \VUgras{patung dada}
\textsuperscript{(59)} menurut alam.

%%K t01-16-1 p
16. --- \VUgras{Ruang musik} tampak sangat menarik. Di sana orang
belajar membaca \VUgras{musik,} menyanyikan \VUgras{not-not}
\textsuperscript{(67)} yang ditulis di atas \VUgras{garis paranada}
(do, re, mi, fa, sol, la, si\,=\,C. D. E. F. G. A. B.), mengenal
\VUgras{kunci-kunci} dan menyanyikan \VUgras{lagu-lagu} yang merdu.
Lembaran musik diletakkan di \VUgras{standar not}
\textsuperscript{(65)}. Salah seorang murid \VUgras{bermain piano}
\textsuperscript{(64)} dan \VUgras{guru musik} \textsuperscript{(66)}
\VUgras{mengetuk irama} \textsuperscript{(68)}.

%%K t01-17-1 p
17. --- Kita mulai melihat sudut \VUgras{bengkel} untuk
\VUgras{kerajinan tangan} \textsuperscript{(69)}, tempat anak-anak
boleh belajar menyerut kayu dan mengikir besi. Hal itu tidak ada di
semua sekolah menengah.

%%K t01-tit-7 sub
\VUsaut{5.0mm}
\VUpk{10.2pt}{40}{Ruang ilmu pengetahuan.}
\VUsaut{3.6mm}

%%K t01-18-1 p
18. --- Guru matematika mengajarkan \VUgras{kepada} murid-murid
\VUgras{ilmu ukur, ilmu hitung, berhitung} dan \VUgras{sistem
metrik}. Di papan tulis kita melihat bangun-bangun geometri yang
utama : \VUgras{lingkaran} \textsuperscript{(a)}, \VUgras{bujur
sangkar} \textsuperscript{(b)}, \VUgras{segitiga}
\textsuperscript{(c)}, \VUgras{garis} \textsuperscript{(d)}.

%%K t01-18-2 p
Kita juga melihat sebuah \VUgras{operasi hitung} ; itu bukan
\VUgras{penjumlahan,} bukan \VUgras{pengurangan,} bukan pula
\VUgras{perkalian} — itu memang \VUgras{pembagian}
\textsuperscript{(e)}.

%%K t01-19-1 p
19. --- Pada papan sistem metrik kita melihat : \VUgras{meter}
\textsuperscript{(a)}, \VUgras{timbangan} \textsuperscript{(b)} dengan
\VUgras{anak-anak timbangannya,} \VUgras{liter} \textsuperscript{(e)},
\VUgras{dekaliter} \textsuperscript{(f)}, \VUgras{desiliter} dan
\VUgras{meter kubik} \textsuperscript{(d)}.

%%K t01-19-2 p
Murid-murid juga mempelajari \VUgras{ilmu alam} \textsuperscript{(11)},
ilmu hewan \textsuperscript{(a)}, ilmu tumbuhan \textsuperscript{(b)},
ilmu bumi \textsuperscript{(c)} dan \VUgras{sejarah}
\textsuperscript{(12)}.

%%K t01-19-3 p
Di dalam lemari ada alat-alat untuk \VUgras{fisika}
\textsuperscript{(15)} dan untuk \VUgras{kimia} \textsuperscript{(16)}.

%%K t01-19-4 p
Di atas lemari ada : \VUgras{bola} \textsuperscript{(14)},
\VUgras{kerucut} dan \VUgras{bola dunia} \textsuperscript{(13)}.

%%K t01-19-5 p
Di atas papan-papan itu tergantung sebuah \VUgras{peta} yang
menggambarkan kelima bagian dunia : \VUgras{Eropa}
\textsuperscript{(g)}, \VUgras{Asia} \textsuperscript{(e)},
\VUgras{Afrika} \textsuperscript{(f)}, \VUgras{Amerika}
\textsuperscript{(ab)} dan \VUgras{Oseania} \textsuperscript{(h)},
serta kedua samudra besar, \VUgras{Pasifik} \textsuperscript{(c)} dan
\VUgras{Atlantik} \textsuperscript{(d)}.

\VUsaut{6.0mm}
\VUfilet{22mm}
